\defmodule{TESTLAT}

This module offers a procedure to test several projections of the lattice specified by a modulus 
and  a set of generating vectors.  Several categories of projections can be considered, with 
specific projections in each category.  All these parameters are specified via 
{\tt InitTestConfig}, stored in a {\tt TestConfig} structure which is passed to the
{\tt TestLattice} procedure.


\bigskip\hrule
\code

DEFINITION MODULE TESTLAT;

FROM BASIS    IMPORT BVect;
FROM MULT     IMPORT MScal, MVect;
FROM CONFIG   IMPORT MDim, BDim, DimProj, MCoord, Categ, OutputType;
FROM NORM     IMPORT CriterionType, Normaliz, NormType, GammaType;
FROM LATBASIS IMPORT Basis;
FROM REDBAS   IMPORT VerifyType;
IMPORT BASISLR;


TYPE
   CoordSet = SET OF MCoord;

   TestConfig = POINTER TO InfoTestConfig;
   InfoTestConfig = RECORD
      m             : MScal;
\endcode
\tabbb
  The modulus.
\endtabbb
\code

      k             : MDim;
\endcode
\tabbb
  Number of generating vectors (usually equal to the order of the
  linear recurrence).
\endtabbb
\code

      NbCoord       : MCoord;
\endcode
\tabbb
  Total number of coordinates that are considered.
\endtabbb
\code

      Maxds         : BDim;
      Maxdp         : DimProj;
\endcode
\tabbb
  {\tt Maxds} is the maximum dimension of the projections over successive 
  coordinates.
  {\tt Maxdp} is the maximum dimension of the projections over nonsuccessive
  coordinates.
  These are the values used for memory allocation in {\tt InitTestConfig}.
\endtabbb
\code

      NbCats, C     : Categ;
\endcode
\tabbb
  Number of categories (or lists) that are considered.
  The variable {\tt C} is initially set to {\tt NbCats} by {\tt TestLattice},
  and is decreased when the generator is not good enough for the 
  category currently considered.
\endtabbb
\code

      DimenProj     : ARRAY Categ OF BDim;
      RangeProj     : ARRAY Categ OF ARRAY DimProj OF BDim;
      DimRegLag     : ARRAY Categ OF BDim;
      RangeRegLag   : ARRAY Categ OF ARRAY DimProj OF BDim;
\endcode
\tabbb
   These arrays specify which projections are considered for each category.
   {\tt DimenProj[c]} is the maximum number of dimensions for which the
   projections over non-successive and non-regular indices are considered.
   (E.g., if the value is 3, the pairs and triples are considered.)
   For $j>1$, {\tt RangeProj[c,j]} indicates the maximal value of $i_j+1$
   that is considered for $j$-dimensional projections in category $c$,
   assuming that $i_1=0$.  For $j=1$, it gives the maximum dimension for
   the successive dimensions, for category $c$.
   Non-successive indices that are regularly spaced, i.e., where
   $i_j = (j-1)s$ for some $s$, are also considered.
   For category $c$, {\tt DimRegLag[c]} gives the maximum value of $s$
   and {\tt RangeRegLag[c,s]} gives the maximum value of $j$ (number
   of dimensions) that is considered for this $s$.
\endtabbb
\code

      Mins2, Maxs2  : ARRAY Categ OF LONGREAL;
\endcode
\tabbb
   For each category $c$, gives the minimum and maximum value of the 
   worst-case criterion within category $c$ for a generator to be 
   considered for category $c$ and the categories above it.
   Whenever {\tt ValCrit[c]} is outside the range
   [{\tt Mins2[c], Maxs2[c]}], the corresponding generator need not be
   considered any longer for all the categories $c'\ge c$.
   It may still be considered, however, for the categories $c'<c$.
\endtabbb
\code

      InvertL, InvertS : BOOLEAN;
\endcode
\tabbb
   Indicates if we want to print the multiplicative inverses of
   the length of the shortest vector and/or the figure of merit {\tt S2}.
   For the spectral test with the Euclidean norm, for instance, 
   if {\tt InvertL = TRUE}, we obtain the distance between hyperplanes
   instead of the length of the shortest vector.
   Default values are {\tt FALSE}.
\endtabbb
\code

      Criterion     : CriterionType;
\endcode
\tabbb
   Indicates which criterion is used. 
\endtabbb
\code

      Normal        : Normaliz;
\endcode
\tabbb
   {\tt Normal} specifies the normalization method
   used to compute {\tt S2}.
\endtabbb
\code

      Verify        : VerifyType;
\endcode
\tabbb
   Indicates the level of verification (i.e., {\tt REDBAS2} or
   {\tt REDBAS3} is used).
\endtabbb
\code

      Output        : OutputType;
\endcode
\tabbb
   Indicates what kind of output is written during the test.
\endtabbb
\code

      MasterGen     : ARRAY MDim OF ARRAY MCoord OF MScal;
\endcode
\tabbb
  {\tt MasterGen} contains an initial set of $k$ generating vectors
  for all the {\tt NbCoord} coordinates that are to be considered in the 
  projections.
  The projections of these vectors, obtained by selecting the 
  appropriate columns in {\tt MasterGen}, together with additional vectors
  of the form $(\dots, m, \dots)$, are used to build the triangular
  bases in {\tt TBas}.
\endtabbb
\code

      TBas          : Basis;
\endcode
\tabbb
  Contains a triangular version of the basis
  for the {\tt d} coordinates that are currently under consideration.
  Used by the {\tt TestProjections} procedure.
\endtabbb
\code

      RBas          : ARRAY DimProj OF Basis;
\endcode
\tabbb
  {\tt RBas[d]} contains a reduced basis
  for the {\tt d} coordinates that are currently under consideration.
  This is used by the {\tt TestProjections} procedure, which goes
  back and forth in this array while exploring the tree of projections.
  {\tt RBas[0]} is used for the lattices that correspond to successive
  coordinates.
  
\endtabbb
\code

      L2, S2        : ARRAY DimProj OF BASISLR.BVect;
\endcode
\tabbb
   The vector {\tt L2} will contain the lengths of the shortest vectors 
   raised to the power $p$ in case of the spectral test with the 
   $\cL_p$ norm.  More precisely, {\tt L2[1,t]} is the value for the
   lattice that corresponds to successive coordinates in $t$ dimensions,
   whereas {\tt L2[d,t]}, for $d>1$, is the worst value for all the 
   $d$-dimensional projections over coordinate sets $I = \{i_1,\dots,i_d\}$
   where $i_1=0$ and $i_d=t-1$.  In all cases,   
   {\tt S2[d,t]} will contain the corresponding normalized quantity,
   also raised to the power $p$, or the Beyer quotient in case of the 
   Beyer test.
\endtabbb
\code

      ValCrit : ARRAY Categ OF LONGREAL;
\endcode
\tabbb
   {\tt ValCrit[c]} will contain the worst-case value of the figure 
   of merit in {\tt S2[d,t]} over the set of projections that are considered
   for category {\tt c}.
\endtabbb
\code

      OK : BOOLEAN;
\endcode
\tabb
   Indicates if the last operation (e.g., call to {\tt TestLattice,
   SpectralTest}, etc.) has been successful.
\endtabb
\code
      END;



PROCEDURE CreateTestConfig (VAR T : TestConfig);
\endcode
  \tab  
  Creates the structure {\tt T}.  One must then call
  {\tt InitTestConfig} to allocate memory for the elements of
  {\tt MasterGen, TBas, RBas, L2}, and {\tt S2}.
  \endtab
\code

PROCEDURE DeleteTestConfig (VAR T : TestConfig);
\endcode
  \tab  
  Deletes the structure {\tt T}, after deallocating the memory
  allocated to its fields by {\tt InitTestConfig}.
  \endtab
\code

PROCEDURE InitTestConfig 
  ( T : TestConfig;  Norm : NormType;
    m  : MScal;   k  : MDim;      NbCoord : MCoord;  
    ds : BDim;    dp : DimProj;   NbCats  : Categ;
    gamma : GammaType;   beta : LONGREAL
  );
\endcode
  \tab
  Initializes the structure {\tt T}.
  In the parameters, {\tt m} represent the modulus, {\tt k} the number of
  generating vectors, 
  {\tt NbCoord} the total number of coordinates to be considered,
  {\tt ds} the maximum dimension of the projections over successive 
  coordinates,
  {\tt dp} the maximum dimension of the projections over nonsuccessive
  coordinates,
  and {\tt NbCats} the number of categories.
  This procedure creates the elements of {\tt MasterGen[1..k, 1..NbCoord]},
  {\tt L2}, and {\tt S2},
  creates the bases {\tt TBas} and {\tt RBas[0..d]},
  initializes their sizes appropriately, their modulus to $m$,
  and their {\tt Norm} field to {\tt Norm}.
  The parameters {\tt gamma {\rm and} beta} are used to initialize {\tt
  T\^{}.Normal} by calling {\tt InitNormaliz}.
  \endtab
\code

PROCEDURE AddToBasis (T : TestConfig;  d : BDim;  Ind : CoordSet;  n : BDim);
\endcode
  \tab
  Assumes that {\tt MasterGen} has been initialized to a proper set
  of generating vectors.
  Adds a vector to the basis {\tt TBas} that corresponds to adding 
  the $j$th coordinate of the generating vectors, and updates the dual.
  \endtab
\code

PROCEDURE RemoveFromBasis (T : TestConfig; d : BDim;  Ind : CoordSet; n : BDim);
\endcode
  \tab
  Removes from the basis {\tt TBas} the vector that corresponds to the 
  $j$th coordinate, and updates the dual.
  \endtab
\code

PROCEDURE TestLattice (T : TestConfig);
\endcode
 \tab  Applies the tests specified by {\tt T} to the lattice
   whose set of generating vectors is in its field \verb=MasterGen=.
 \endtab
\code

END TESTLAT.
\endcode
